\documentclass{../note}

\begin{document}

\begin{zadanie}{1}
Niech $a, b, c>0$. Określamy ciągi $(a_n)$, $(b_n)$ i $(c_n)$ rekurencyjnie, $a_{0}=a$, $b_{0}=b$, $c_{0}=c$,
\[ a_{n+1}=\frac{3}{\frac{1}{a_{n}}+\frac{1}{b_{n}}+\frac{1}{c_{n}}}, \quad b_{n+1}=\sqrt[3]{a_{n}b_{n}c_{n}}, \quad c_{n+1}=\frac{a_{n}+b_{n}+c_{n}}{3}, \quad n \ge 0. \]
Wykaż, że te trzy ciągi są zbieżne do tej samej granicy skończonej.
\end{zadanie}
\begin{rozwiazanie}
Z indukcji matematycznej możemy wywnioskować, że $a_n \leq b_n \leq c_n$ dla $n \geq 1$, bo zarówno baza i krok indukcyjny wynikają z nierówności pomiędzy średnimi. Wszystkie średnie są pomiędzy minimalnym a maksymalnym argumentem, czyli $a_n \leq a_{n+1}$ oraz $c_{n+1} \leq c_n$, czyli ciągi $a_n$ i $c_n$ są odpowiednio rosnące i malejące. Oba są ograniczone, odpowiednio od góry przez $c$ i od dołu przez $a$, więc oba mają granicę. Załóżmy nie wprost, że $\lim_{n\rightarrow\infty} a_n = A, \lim_{n\rightarrow\infty} c_n = C$ oraz że $A < C$. Wtedy dla każdego $\varepsilon > 0$ i ddd $n$ zachodzi $c_n - C < \varepsilon$, jednak wtedy 
\[c_{n+1} = \frac{a_{n}+b_{n}+c_{n}}{3} < \frac{A+2(C + \varepsilon)}{3} = \frac{A-C+2\varepsilon}{3}+ C\]
czyli $c_{n+1} < C$ dla $\varepsilon < \frac{C - A}{2}$, co jest sprzecznością. Zatem ciągi $a_n$ i $c_n$ mają tę samą granicę $A = C$ i z tw. o trzech ciągach ciąg $b_n$ też zmierza do tej granicy.
\end{rozwiazanie}

\begin{zadanie}{2}
Obliczyć granice ciągów:
\begin{enumerate}[label=(\alph*)]
    \item $a_{n}=\sqrt[n]{3^{n}-2^{n}}$ \hfill [0,25 p.]
    \item $b_{n}=\sqrt[n]{x^{n}+x^{n-1}+\dots+x+1}$, gdzie $x>0$ jest ustalone. \hfill [0,25 p.]
\end{enumerate}
Dla jakich $x>0$ ciąg z punktu (b) jest monotoniczny? \hfill [0,5 p.]
\end{zadanie}
\begin{rozwiazanie}
\begin{enumerate}[label=(\alph*)]
    \item Mamy $3^{1 - \frac1n} = \sqrt[n]{3^{n-1}} < a_n < \sqrt[n]{3^n} = 3$. Dla $n \rightarrow \infty$ lewa strona zmierza do 3, więc z twierdzenia o trzech ciągach $\lim_{n\rightarrow\infty} a_n = 3$.
    \item Rozważmy najpierw $x < 1$. Wtedy 
    \[1 < b_n = \sqrt[n]{\frac{1-x^{n+1}}{1-x}} < (1-x)^{-\frac{1}{n}}\]
    Prawa strona zbiega do 1, więc z twierdzenia o trzech ciągach $b_n$ też zbiega do $1$. Wiemy też, że $b_0 = 1$ i $b_1 > 1$, więc w tym przypadku ciąg $b_n$ nie jest monotoniczny.\\
    Gdy $x = 1$, to $b_n = \sqrt[n]{n}$. W skrypcie do wykładów jest dowód, że to zbiega do $1$. Z takiego samego argumentu, jak w przypadku $x < 1$, możemy wywnioskować, że $b_n$ nie jest monotoniczny.\\
    Gdy $x > 1$ zachodzi 
    \[b_n = \frac{x\sqrt[n]{x-\frac{1}{x^n}}}{\sqrt[n]{x-1}}\]
    $\sqrt[n]{x-\frac{1}{x^n}}$ jest ograniczone od dołu przez 1 i od góry przez $\sqrt[n]{x}$, czyli zbiega do 1. $\sqrt[n]{x-1}$ też zbiega do 1, więc $b_n$ zbiega do $x$. Podobnie jak w pozostałych przypadkach, $b_0 = 0$, $b_1 = x+1 > 1$, a ddd $n$ zachodzi $b_n < x + 1$, czyli ciąg nie jest monotoniczny.
\end{enumerate}
\end{rozwiazanie}
\begin{zadanie}{3}
Niech $(a_{n})_{n\ge1}$, $(b_{n})_{n\ge1}$ będą ciągami liczb rzeczywistych. Przypuśćmy, że $\lim_{n\rightarrow\infty}a_{n}=A$ i $\lim_{n\rightarrow\infty}b_{n}=B$. Oblicz granicę ciągu
\[ c_{n}=\frac{a_{1}b_{n}+a_{2}b_{n-1}+\dots+a_{n}b_{1}}{n} \]
\end{zadanie}
\begin{rozwiazanie}
Obliczmy granicę dla parzystych $n$. Niech $n = 2m$. Oznaczmy $\alpha_n = a_n - A, \beta_n = b_n - B$. Zachodzi:
\[|c_n - AB| = \left|\frac1{2m}\left(\sum_{k=1}^{m}a_{m+k}b_{m-k} - AB\right) + \frac1{2m}\left(\sum_{k=1}^{m}|a_{m-k}b_{m+k} - AB|\right)\right| =\]
\[\left|\frac1{2m}\left(\sum_{k=1}^{m}\alpha_{m+k}b_{m-k} + A\beta_{m-k}\right) + \frac1{2m}\left(\sum_{k=1}^{m}\beta_{m+k}a_{m-k} + B\alpha_{m-k}\right)\right|\]
Z nierówności trójkąta to jest mniejsze lub równe
\begin{equation}
\frac{1}{2m}\left(\sum_{k=1}^{m}|\alpha_{m+k}b_{m-k}| + |\beta_{m+k}a_{m-k}|\right) + \frac{A}{2m}\left|\sum_{k=1}^{m}\beta_k\right| + \frac{B}{2m}\left|\sum_{k=1}^{m}\alpha_k\right|
\end{equation}
Niech $a_{sup}$ i $b_{sup}$ to suprema danych ciągów. Dla dowolnego $\varepsilon > 0$ i ddd $m$, pierwszy składnik (1) jest ograniczony od góry przez
\[\frac{1}{2m}\left(\sum_{k=1}^{m}|\varepsilon b_{sup}| + |\varepsilon a_{sup}|\right) = \frac{\varepsilon}{2}(a_{sup} + b_{sup})\]
Z warunku Cauchy'ego istnieje takie $l$, że dla każdego $m$ zachodzi $\sum_{k=l}^{m}\beta_k < \varepsilon$. $\sum_{k=1}^{l}\beta_k$ jest stałe, więc ddd $m$ zachodzi $\frac1{2m}\left|\sum_{k=1}^{m}\beta_k\right| < 2\varepsilon$. Zatem (1) jest ograniczone od góry przez 
\[\frac{9\varepsilon}{2}(a_{sup} + b_{sup})\]
Zatem $|c_n - AB|$ zbiega do 0, czyli $AB$ jest granicą ciagu $c_n$. Dla nieparzystych $n$ można przeprowadzić podobne rozumowanie.
\end{rozwiazanie}
\begin{zadanie}{4}
Obliczyć granice ciągów:
Obliczyć granice ciągów:
\begin{enumerate}[label=(\alph*)]
    \item $a_{n}=\sqrt{4n^{2}+n+1}-\sqrt{n^{2}+n+1}-n,$ \hfill [0,25 p.]
    \item $b_{n}=\frac{n+1}{n^{2}+1}+\frac{n+2}{n^{2}+2}+\dots+\frac{n+n}{n^{2}+n}$ \hfill [0,25 p.]
    \item $\sqrt{c},\sqrt{c+\sqrt{c}},\sqrt{c+\sqrt{c+\sqrt{c}}}, \dots$ w zależności od parametru $c\ge0$. \hfill [0,5 p.] \\
    
\end{enumerate}
\end{zadanie}
\begin{rozwiazanie}
\begin{enumerate}[label=(\alph*)]
    \item Najpierw lemat: Dla dowolnego stałego $a$ zachodzi $\lim_{x\rightarrow\infty}\sqrt{x + a} - \sqrt{x} = 0$. Dowód: $|\sqrt{x + a} - \sqrt{x}| < \sqrt{x + a + \frac{a^2}{4x}} - \sqrt{x} = \frac{a}{2\sqrt{x}}$. Dla ddd $x$ prawa strona zbiega do 0, co kończy dowód lematu. Korzystamy z zasad arytmetyki dla granic:
    \[\lim_{n\rightarrow\infty} a_n = \lim_{n\rightarrow\infty} \sqrt{\left(2n+\frac14\right)^2+\frac{15}{16}}-\sqrt{\left(n+\frac12\right)^2+\frac34}-n=\]
    \[\lim_{n\rightarrow\infty} \sqrt{\left(2n+\frac14\right)^2}-\sqrt{\left(n+\frac12\right)^2}-n = -\frac14\]
    \item Wszystkie mianowniki są większe od $n^2$ i mniejsze lub równe $n(n+1)$. Korzystając z wzoru na sumę pierwszych $n$ liczb naturalnych, mamy:
    \[\frac32 - \frac1n = \frac{1}{n(n+1)}\left(n^2 + \frac{n(n+1)}{2}\right) \leq b_n \leq \frac{1}{n^2}\left(n^2 + \frac{n(n+1)}{2}\right) = \frac32 + \frac{1}{2n}\]
    Zarówno lewa, jak i prawa strona zbiegają do $\frac32$ w granicy, więc z tw. o trzech ciągach $b_n$ też zbiega do $\frac32$.
    \item Niech $g = \frac{1+\sqrt{1+4c}}{2}$ to dodatni pierwiastek wielomianu $x^2=x+c$ i niech $c_n$ ciąg określony wyrażeniem z treści. $c_n = \sqrt{c + c_{n-1}}$, więc zachodzi 
    \[(g-c_n)(g+c_n) = g^2 - c_n^2 = c + g - c - c_{n-1} = g - c_{n-1}\]
    czyli
    \[|g-c_n| = \frac{|g - c_{n-1}|}{g+c_n} < \frac{|g - c_{n-1}|}{g}\]
    $g = \frac{1+\sqrt{1+4c}}{2} > \frac{1+\sqrt{1}}{2} = 1$, czyli $|g-c_n|$ zbiega do 0 w tempie wykładniczym, czyli $g$ jest granicą ciągu $c_n$.
\end{enumerate}
\end{rozwiazanie}
\begin{zadanie}{5}
Niech $a>0$. Zbadaj zbieżność ciągu $(a_{n})$ zadanego rekurencyjnie wzorami
\[a_{0}=a, \hspace{20pt} a_{n+1}=1+\frac{1}{a_{n}}, \hspace{20pt} n\ge0.\] \\
Jeśli ciąg jest zbieżny, wyznacz jego granicę.
\end{zadanie}
\begin{rozwiazanie}
Z definicji ciągu wiemy, że $a_n > 1$ dla $n \geq 1$. Pokażemy, że ciąg zbiega do złotej liczby, czyli dodatniej liczby spełniającej $\phi = 1 + \frac1\phi$. Zachodzi:
\[|\phi - a_{n+1}| = \left|1 + \frac1\phi - 1 - \frac1{a_n}\right| = \frac{|\phi - a_n|}{\phi a_n} < \frac{|\phi - a_n|}{\phi a_n}\]
Zatem $|\phi - a_n|$ maleje wykładniczo, czyli zbiega do 0, czyli $a_n$ zbiega do $\phi$.
\end{rozwiazanie}
\begin{zadanie}{6}
Niech $f:[0,1]\rightarrow[0,1]$ będzie funkcją niemalejącą. Czy musi istnieć $a\in[0,1]$, dla którego $f(a)=a$?
\end{zadanie}
\begin{rozwiazanie}
Niech $X = \left\{a : f(a) > a\right\}$ i $x = \sup X$. Jeśli $X$ jest pusty, to $f(0) \leq 0$, czyli $f(0) = 0$. Jeśli $x \in X$, to 
\[f\left(\frac{f(x)+x}{2}\right) \geq f(x) > \frac{f(x)+x}{2}\]
co jest sprzeczne z założeniem, że $x = \sup X$, bo $\frac{f(x)+x}{2} > x$. Musi więc istnieć taki ciąg $(x_n)_{n\geq1} \subseteq X$, którego granicą jest $x$. Jeśli $f(x) < x$, to $f(x) < \frac{f(x)+x}{2} < x$, czyli dla ddd $n$ zachodzi 
\[f(x_n) > x_n > \frac{f(x)+x}{2} > f(x) \geq f(x_n)\]
co jest sprzecznością. Zatem musi zachodzić $f(x) = x$, czyli zawsze istnieje takie $a \in [0, 1]$, że $f(a) = a$.
\end{rozwiazanie}
\end{document}